\section{Fundamentos, potência, energia e demanda}

\begin{frame}{Grandezas em corrente alternada}
\small
\begin{columns}[T,onlytextwidth]
\column{0.54\textwidth}
\begin{align*}
v(t)&=V_m\sen(\omega t+\theta_v)\\
i(t)&=I_m\sen(\omega t+\theta_i)\\
V&=\frac{V_m}{\sqrt{2}},\qquad I=\frac{I_m}{\sqrt{2}}\\
\varphi&=\theta_v-\theta_i
\end{align*}
\begin{itemize}
\item RMS relaciona a grandeza CA ao efeito térmico equivalente em CC.
\item Em carga indutiva, corrente tende a atrasar a tensão.
\item Em carga capacitiva, corrente tende a adiantar a tensão.
\end{itemize}
\column{0.44\textwidth}
\begin{tikzpicture}
\begin{axis}[width=5.5cm,height=3.7cm,axis lines=middle,xlabel={$\omega t$},xtick={0,90,180,270,360},xticklabels={$0$,$90$,$180$,$270$,$360$},ytick=\empty,domain=0:360,samples=120,ymin=-1.3,ymax=1.3,grid=both]
\addplot[corPrim,thick] {sin(x)};
\addplot[corAccent,thick,dashed] {0.85*sin(x-35)};
\end{axis}
\node[corPrim,font=\scriptsize] at (1.25,3.25) {$v(t)$};
\node[corAccent,font=\scriptsize] at (4.25,2.45) {$i(t)$};
\end{tikzpicture}
\end{columns}
\end{frame}

\begin{frame}{Potência ativa, reativa e aparente}
\small
\begin{columns}[T,onlytextwidth]
\column{0.50\textwidth}
\begin{align*}
S &= VI \quad [\si{VA}]\\
P &= VI\cos\varphi \quad [\si{W}]\\
Q &= VI\sen\varphi \quad [\si{var}]\\
S^2 &= P^2+Q^2,\qquad fp=\frac{P}{S}
\end{align*}
\begin{itemize}
\item $P$: potência convertida em trabalho, calor, luz etc.
\item $Q$: intercâmbio de energia reativa.
\item $S$: capacidade elétrica exigida de cabos, quadros e transformadores.
\end{itemize}
\column{0.48\textwidth}
\centering
\begin{tikzpicture}[scale=0.9]
\draw[->,thick] (0,0)--(4.1,0) node[right]{$P$};
\draw[->,thick] (4.1,0)--(4.1,2.4) node[above]{$Q$};
\draw[->,thick,corPrim] (0,0)--(4.1,2.4) node[midway,above left]{$S$};
\draw (0.8,0) arc (0:30:0.8);
\node at (1.15,0.28) {$\varphi$};
\end{tikzpicture}
\end{columns}
\end{frame}

\begin{frame}{Sistemas trifásicos}
\small
\begin{columns}[T,onlytextwidth]
\column{0.52\textwidth}
\begin{align*}
P_{3\phi}&=\sqrt3\,V_{LL}I_L fp\\
Q_{3\phi}&=\sqrt3\,V_{LL}I_L\sen\varphi\\
S_{3\phi}&=\sqrt3\,V_{LL}I_L
\end{align*}
\begin{itemize}
\item Cargas equilibradas reduzem corrente de neutro fundamental.
\item Cargas monofásicas devem ser distribuídas entre as fases.
\item Harmônicas triplas podem somar no neutro mesmo com correntes eficazes semelhantes nas fases.
\end{itemize}
\column{0.46\textwidth}
\centering
\begin{tikzpicture}[scale=0.95]
\draw[->,thick,corPrim] (0,0)--(2,0) node[right]{$V_A$};
\draw[->,thick,corDef] (0,0)--({2*cos(120)},{2*sin(120)}) node[left]{$V_B$};
\draw[->,thick,corAccent] (0,0)--({2*cos(240)},{2*sin(240)}) node[left]{$V_C$};
\draw[thin,black!25] (0,0) circle (2);
\node at (0,-2.35) {defasagem de $120^\circ$};
\end{tikzpicture}
\end{columns}
\end{frame}

\begin{frame}{Corrente de projeto}
\small
\begin{columns}[T,onlytextwidth]
\column{0.53\textwidth}
\begin{caixaDef}
A corrente de projeto $I_b$ é a corrente prevista no circuito considerando potência, tensão, rendimento, fator de potência e regime de funcionamento.
\end{caixaDef}
\begin{align*}
I_{1\phi}&=\frac{P}{V\,fp\,\eta}\\
I_{3\phi}&=\frac{P}{\sqrt3\,V_{LL}\,fp\,\eta}
\end{align*}
\column{0.45\textwidth}
\begin{caixaEx}
Motor de $\SI{15}{kW}$, $\SI{380}{V}$, $fp=0{,}85$, $\eta=0{,}92$:
\[
I_b=\frac{15000}{\sqrt3\cdot380\cdot0{,}85\cdot0{,}92}\approx\SI{29.2}{A}.
\]
Esse valor ainda não resolve partida, cabo, disjuntor ou queda de tensão.
\end{caixaEx}
\end{columns}
\end{frame}

\begin{frame}{Carga instalada, demanda e diversidade}
\small
\begin{columns}[T,onlytextwidth]
\column{0.50\textwidth}
\begin{align*}
P_{inst}&=\sum P_k\\
D_{max}&=\max_t P(t)\\
F_d&=\frac{D_{max}}{P_{inst}}\\
F_{div}&=\frac{\sum D_{max,k}}{D_{max,grupo}}
\end{align*}
\begin{itemize}
\item A entrada não deve ser dimensionada supondo simultaneidade total sem justificativa.
\item Fator de demanda deve vir de regra aplicável, histórico, processo ou critério documentado.
\end{itemize}
\column{0.48\textwidth}
\begin{tikzpicture}
\begin{axis}[width=5.7cm,height=3.6cm,xlabel={hora},ylabel={$P$ (kW)},xmin=0,xmax=24,ymin=0,ymax=100,grid=both,xtick={0,6,12,18,24}]
\addplot[corPrim,thick,smooth] coordinates {(0,18)(3,15)(6,30)(9,54)(12,71)(15,63)(18,90)(21,50)(24,26)};
\addplot[corAccent,dashed,thick] coordinates {(0,90)(24,90)};
\end{axis}
\end{tikzpicture}
\end{columns}
\end{frame}

\begin{frame}{Levantamento de cargas moderno}
\scriptsize
\begin{tabularx}{\textwidth}{>{\bfseries}p{2.6cm}p{3.1cm}X}
\toprule
Carga & Dados mínimos & O que acrescentar ao projeto \\
\midrule
Iluminação & Potência, comando, ambiente & Iluminância, emergência, automação e eficiência. \\
TUG/TUE & Potência, tensão, uso & Circuitos, DR, simultaneidade, expansão. \\
Motores & kW, regime, partida & $I_N$, inrush, torque, proteção, VFD/soft-starter. \\
TI/eletrônica & W/VA, UPS & Harmônicas, DPS, continuidade, aterramento funcional. \\
Ar-condicionado & Potência e tecnologia & Inverter, corrente máxima, diversidade e partida. \\
VE & Potência de recarga, modo & Circuito dedicado, DR, gestão de carga, expansão. \\
FV/armazenamento & Potência, inversor & Proteções CA/CC, anti-ilhamento, fluxos bidirecionais. \\
Processo industrial & Ciclo e criticidade & Continuidade, redundância, seletividade, produção. \\
\bottomrule
\end{tabularx}
\end{frame}

%=====================================================================
